\documentclass[11pt,a4paper]{article}

% ------------------------------------------------------------------
% LogicWard — Demo Guide (presenter runbook)
% Compile:  pdflatex (run twice)  or  latexmk -pdf
% ------------------------------------------------------------------
\usepackage[utf8]{inputenc}
\usepackage[T1]{fontenc}
\usepackage{lmodern}
\usepackage[margin=2.1cm]{geometry}
\usepackage[dvipsnames,table]{xcolor}
\usepackage{booktabs}
\usepackage{tabularx}
\usepackage{array}
\usepackage{listings}
\usepackage{titlesec}
\usepackage{enumitem}
\usepackage{fancyhdr}
\usepackage{parskip}
\usepackage[most]{tcolorbox}
\usepackage{pdflscape}
\usepackage{underscore}
\usepackage[hidelinks]{hyperref}

\definecolor{navy}{HTML}{14304F}
\definecolor{navy2}{HTML}{1D4370}
\definecolor{accent}{HTML}{3987E5}
\definecolor{ink}{HTML}{1F2937}
\definecolor{muted}{HTML}{5B6B80}
\definecolor{good}{HTML}{0E7A0E}
\definecolor{warn}{HTML}{A16207}
\definecolor{serious}{HTML}{C05621}
\definecolor{critical}{HTML}{C02626}
\definecolor{page}{HTML}{EEF0F4}
\definecolor{hair}{HTML}{DBE0E8}
\definecolor{codebg}{HTML}{F6F8FB}
\definecolor{critbg}{HTML}{FBE8E8}
\definecolor{highbg}{HTML}{FBEEE6}
\definecolor{medbg}{HTML}{FDF4E3}
\definecolor{lowbg}{HTML}{EEF1F5}

\hypersetup{colorlinks=true,linkcolor=navy2,urlcolor=accent}

\titleformat{\section}{\Large\bfseries\color{navy}}{\thesection}{0.6em}{}
\titleformat{\subsection}{\large\bfseries\color{navy2}}{\thesubsection}{0.6em}{}
\titlespacing*{\section}{0pt}{1.3em}{0.5em}

\pagestyle{fancy}\fancyhf{}
\setlength{\headheight}{24pt}\addtolength{\topmargin}{-10pt}
\renewcommand{\headrulewidth}{0.4pt}
\fancyhead[L]{\small\color{muted}\textbf{LogicWard}\,\textbullet\, Demo Guide}
\fancyhead[R]{\small\color{muted}Presenter Runbook}
\fancyfoot[C]{\small\color{muted}\thepage}

\lstdefinestyle{lw}{basicstyle=\ttfamily\small\color{ink},backgroundcolor=\color{codebg},
  frame=single,framesep=6pt,rulecolor=\color{hair},breaklines=true,columns=fullflexible,
  keepspaces=true,showstringspaces=false,xleftmargin=4pt,xrightmargin=4pt,aboveskip=7pt,belowskip=7pt}
\lstset{style=lw}

\newtcolorbox{say}[1][]{colback=accent!7,colframe=accent,boxrule=0.8pt,left=8pt,right=8pt,
  top=5pt,bottom=5pt,arc=2pt,#1}
\newtcolorbox{note}[1][]{colback=page,colframe=hair,boxrule=1pt,left=8pt,right=8pt,
  top=5pt,bottom=5pt,arc=2pt,#1}

\newcommand{\code}[1]{\texttt{\small #1}}
\newcolumntype{L}{>{\raggedright\arraybackslash}X}
\newcommand{\sev}[2]{\colorbox{#1}{\small\strut\,#2\,}}
\newcommand{\crit}{\sev{critbg}{\color{critical}\textbf{critical}}}
\newcommand{\high}{\sev{highbg}{\color{serious}\textbf{high}}}
\newcommand{\med}{\sev{medbg}{\color{warn}\textbf{medium}}}
\newcommand{\low}{\sev{lowbg}{\color{muted}\textbf{low}}}

\begin{document}

% ---------------- title ----------------
\begin{titlepage}
\thispagestyle{empty}
\vspace*{1.4cm}
{\color{navy}\rule{\textwidth}{2.5pt}}\\[0.5cm]
{\Huge\bfseries\color{navy} LogicWard --- Demo Guide}\\[0.35cm]
{\LARGE\color{navy2} How to run the demo, and read every detection}\\[0.5cm]
{\large\color{ink} A step-by-step presenter runbook: which script to run from where,
what the attacker does, where it lands, how LogicWard detects it, which dashboard view
shows it, and the exact MITRE ATT\&CK for ICS technique behind each alert.}\\[0.4cm]
{\color{navy}\rule{\textwidth}{2.5pt}}\\[1.0cm]

\begin{note}
\large
\textbf{The whole demo is one command:}\\[4pt]
\code{python -m logicward.attacker.demo_sequence}\\[6pt]
Open \textbf{http://localhost:8080/} (login \code{soc/soc123}) and watch the attacks land.
A full run produces \textbf{14 events --- 4 critical, 2 high, 3 medium, 5 low}.
\end{note}

\vfill
{\small\color{muted} Repo: https://github.com/positromen/Adani-Project-OT \quad\textbullet\quad
See also USAGE.md (full how-to) and DESIGN.md / LogicWard\_Design.pdf (architecture).}
\end{titlepage}

\tableofcontents
\newpage

% ==================================================================
\section{The cast --- who is who, and where things run}

LogicWard has three roles. In the \textbf{one-command demo} all three run on your
laptop (an \emph{embedded} PLC); in the \textbf{real deployment} they are three separate
machines on the same Wi-Fi. The detection logic is identical either way.

\begin{tabularx}{\textwidth}{@{}p{3.1cm} L p{4.4cm}@{}}
\toprule
\textbf{Role} & \textbf{What it is} & \textbf{Where it runs} \\
\midrule
\textbf{The PLC (victim)} & Thermal power-plant controller: a Modbus TCP server + its
  ladder program (a real Rockwell \code{.L5X}) + program-download endpoint. &
  Raspberry Pi (real) \emph{or} embedded in the demo process (laptop). \\
\addlinespace
\textbf{The Attacker} & A toolkit that tampers with the PLC: setpoint writes, program
  downloads, a Modbus flood, a rogue device. & \code{logicward.attacker} --- a 2nd box
  (real) \emph{or} the demo script (laptop). \\
\addlinespace
\textbf{LogicWard (defender)} & The detection engine + SOC dashboard that diffs the live
  PLC against a signed baseline and raises alerts. & The laptop --- \code{localhost:8080}. \\
\bottomrule
\end{tabularx}

\begin{note}
\textbf{The mental model for the audience:} \emph{``An attacker who can reach a power-plant
PLC quietly rewrites its safety logic. LogicWard sits beside the PLC, holds a
cryptographically-signed copy of the approved program, and the instant the running logic
or a live value drifts from that baseline it raises a severity-ranked, MITRE-mapped
alert --- naming what changed, and how.''}
\end{note}

% ==================================================================
\section{Running the demo}

\subsection{Option A --- one command (recommended for judging)}
From the project root on the laptop:
\begin{lstlisting}
python -m logicward.attacker.demo_sequence          # ~90s, watchable pacing
python -m logicward.attacker.demo_sequence --fast    # ~15s, quick run
\end{lstlisting}
This brings up the SOC dashboard + an embedded thermal PLC + the program endpoints, then
runs the whole attack story with narration in the terminal. Open
\textbf{http://localhost:8080/} (login \code{soc/soc123}) and keep it on screen.

\subsection{Option B --- drive it yourself}
\begin{lstlisting}
# Terminal 1 - the SOC dashboard (embedded PLC)
python -m logicward.dashboard.app

# Terminal 2 - the program endpoints (so program-download attacks land)
python -m logicward.plant.logic_store

# Terminal 3 - fire attacks one at a time (see the catalogue in section 4)
python -m logicward.attacker.attacks logic-inversion
python -m logicward.attacker.attacks ddos --count 800
\end{lstlisting}

\subsection{Option C --- the real two-machine setup}
Pi runs the PLC + sensor agent; the laptop runs the dashboard; a third box attacks.
Full runbook in \textbf{USAGE.md \S7}. In short:
\begin{lstlisting}
# On the Pi:      bash deploy/pi_bootstrap.sh <laptop-ip>   &&   bash deploy/run_pi.sh
# On the laptop:  .\deploy\run_laptop.ps1 -PiHost siddhesh-pi.local
# From attacker:  python -m logicward.attacker.attacks --host siddhesh-pi.local logic-inversion
\end{lstlisting}

\begin{say}
\textbf{Say:} ``Everything you're about to see is two computers talking over Wi-Fi --- an
attacker on one, the plant and its guardian on the other. For this screen it's one laptop,
but the code path is identical.''
\end{say}

% ==================================================================
\section{The story, act by act}

Each act below lists: the \textbf{command}, \textbf{where the attack lands}, \textbf{how
LogicWard detects it}, \textbf{which dashboard view shows it}, the \textbf{MITRE ATT\&CK
for ICS} technique, and a \textbf{line to say}.

\subsection{Act 1 --- Normal operations}
Nothing is attacked. The plant runs nominal, the baseline is HMAC-locked, the live program
matches the baseline, the alert feed is empty.
\begin{itemize}[leftmargin=1.3em,itemsep=1pt]
  \item \textbf{Dashboard:} \emph{Overview} --- baseline integrity \code{VALID}, ``program
        in sync'', zero alerts. \emph{Live Plant} shows drum level, steam pressure, turbine
        speed all nominal.
\end{itemize}
\begin{say}\textbf{Say:} ``A thermal power plant PLC, continuously verified against a
cryptographically-signed approved program. Quiet board, green baseline.''\end{say}

\subsection{Act 2 --- Reconnaissance: a rogue device}
\textbf{Attack:} an unapproved engineering laptop joins the OT network.
\begin{itemize}[leftmargin=1.3em,itemsep=1pt]
  \item \textbf{Command:} \code{python -m logicward.attacker.attacks rogue}
  \item \textbf{Lands on:} the OT LAN segment (an ARP announcement) --- not the PLC yet.
  \item \textbf{Detected by:} \code{arp\_watch} (physical plane, on the Pi agent) --- MAC not
        on the approved allowlist.
  \item \textbf{Event:} \code{physical.rogue\_device} \quad\med
  \item \textbf{Shown on:} \emph{Alerts} feed.
  \item \textbf{MITRE:} Initial Access \textbullet\ \textbf{Rogue Master} \textbullet\ \code{T0848}
\end{itemize}
\begin{say}\textbf{Say:} ``Before touching the PLC, the attacker's machine appears on the
OT segment. An unknown MAC on this network is itself an alert --- Rogue Master, T0848.''\end{say}

\subsection{Act 3 --- Silent logic tampering}
The core of the demo: five ways an attacker rewrites the PLC's safety logic.

\textbf{(3a) Setpoint drift --- over Modbus.}
\begin{itemize}[leftmargin=1.3em,itemsep=1pt]
  \item \textbf{Command:} \code{python -m logicward.attacker.attacks setpoint-drift} \quad(drum-level trip \code{220 -> 40} mm)
  \item \textbf{Lands on:} a Modbus holding register on the PLC (function code 06 write).
  \item \textbf{Detected by:} the drift engine's \textbf{register diff} (live Modbus vs baseline snapshot).
  \item \textbf{Event:} \code{cyber.setpoint\_drift} \quad\high\ (safety-critical setpoint)
  \item \textbf{Shown on:} \emph{Alerts}, and the changed value on \emph{Program Diff} / \emph{Live Plant}.
  \item \textbf{MITRE:} Impair Process Control \textbullet\ \textbf{Modify Parameter} \textbullet\ \code{T0836}
\end{itemize}

\textbf{(3b--3d) Program downloads --- structural logic changes.} Each command pushes a
tampered \code{.L5X} to the PLC's program-download endpoint; LogicWard canonicalizes and
diffs it, and the passive FIM also fires on the file change.
\begin{itemize}[leftmargin=1.3em,itemsep=2pt]
  \item \textbf{Logic inversion} --- \code{attacks logic-inversion} --- flips the drum-level
        trip compare \code{LES$\rightarrow$GRT} (the interlock now fires backwards).\\
        Event \code{cyber.logic\_inversion} \crit\ \textbullet\ Persistence \textbullet\ \textbf{Modify Program} \textbullet\ \code{T0889}
  \item \textbf{Condition stripping} --- \code{attacks condition-stripping} --- removes the
        running-interlock from the furnace-flame trip (a safety input deleted).\\
        Event \code{cyber.condition\_stripping} \crit\ \textbullet\ Persistence \textbullet\ \textbf{Modify Program} \textbullet\ \code{T0889}
  \item \textbf{Coil hijack} --- \code{attacks coil-hijack} --- repoints the feedwater-trip
        output coil to a different actuator.\\
        Event \code{cyber.coil\_hijack} \crit\ \textbullet\ Persistence \textbullet\ \textbf{Modify Program} \textbullet\ \code{T0889}
\end{itemize}

\textbf{(3e) Rung injection --- a backdoor rung.}
\begin{itemize}[leftmargin=1.3em,itemsep=1pt]
  \item \textbf{Command:} \code{python -m logicward.attacker.attacks rung-injection} (unlatches the turbine trip)
  \item \textbf{Detected by:} structural L5X diff (a new, unapproved rung appears) + FIM.
  \item \textbf{Event:} \code{cyber.rung\_injection} \quad\crit
  \item \textbf{Shown on:} \emph{Program Diff} --- the new rung highlighted green.
  \item \textbf{MITRE:} Lateral Movement \textbullet\ \textbf{Program Download} \textbullet\ \code{T0843}
\end{itemize}
Each program download also raises \code{cyber.program\_file\_modified} \low\ from the
\textbf{FIM} sensor (Persistence \textbullet\ Modify Program \textbullet\ \code{T0889}) --- a second,
independent detection of the same tamper via the file system.

\begin{say}\textbf{Say (on Program Diff):} ``This is the exact logic flip --- the drum-low
trip compare inverted from LES to GRT. A harmless re-export with new timestamps would show
\emph{nothing} here; a single changed instruction lights up red. That's canonicalization
doing its job.''\end{say}

\subsection{Act 4 --- Process manipulation \& impact}
\textbf{(4a) Force a control coil.}
\begin{itemize}[leftmargin=1.3em,itemsep=1pt]
  \item \textbf{Command:} \code{python -m logicward.attacker.attacks force-coil} (forces the fuel-valve coil over Modbus)
  \item \textbf{Detected by:} register diff (a control coil forced away from baseline).
  \item \textbf{Event:} \code{cyber.register\_change} \quad\low
  \item \textbf{MITRE:} Impair Process Control \textbullet\ \textbf{Unauthorized Command Message} \textbullet\ \code{T0855}
\end{itemize}

\textbf{(4b) DDoS flood.}
\begin{itemize}[leftmargin=1.3em,itemsep=1pt]
  \item \textbf{Command:} \code{python -m logicward.attacker.attacks ddos --count 800}
  \item \textbf{Lands on:} the PLC's Modbus server (request flood).
  \item \textbf{Detected by:} the \textbf{resource plane} (CPU spike on the PLC host).
  \item \textbf{Event:} \code{resource.cpu\_spike} \quad\med
  \item \textbf{MITRE:} Inhibit Response Function \textbullet\ \textbf{Denial of Service} \textbullet\ \code{T0814}
\end{itemize}

\subsection{Act 5 --- Physical tamper}
\textbf{Attack:} the Ethernet cable is pulled and the enclosure door is opened.
\begin{itemize}[leftmargin=1.3em,itemsep=1pt]
  \item \textbf{Detected by:} \code{link\_watch} (carrier lost) and \code{gpio\_tamper}
        (door switch) --- physical plane, on the Pi agent.
  \item \textbf{Events:} \code{physical.link\_down} \med\ \textbullet\ \code{physical.enclosure\_open} \high
  \item \textbf{MITRE:} \code{link\_down} $\rightarrow$ Denial of Service \code{T0814};
        enclosure tamper has \emph{no} direct ICS technique (shown honestly as \code{N/A}).
\end{itemize}
\begin{say}\textbf{Say:} ``The same fabric that caught a logic edit also catches a pulled
cable and an opened cabinet --- one event contract across cyber, physical, and resource.''\end{say}

\subsection{Detection summary}
\begin{note}
After a full run the dashboard shows \textbf{14 events}:
\quad\crit\,$\times4$ \quad \high\,$\times2$ \quad \med\,$\times3$ \quad \low\,$\times5$.\\
Every event is in the append-only evidence log --- export the \textbf{signed PDF forensic
report} from the \emph{Evidence} view. \textbf{Say:} ``One click --- a signed incident
report an investigator can trust.''
\end{note}

% ==================================================================
\section{Where each detection shows on the dashboard}
Login \code{soc/soc123}. Roles: \code{operator} (read-only), \code{engineer}
(+baseline), \code{soc} (everything).
\begin{tabularx}{\textwidth}{@{}p{3.2cm} L@{}}
\toprule
\textbf{View} & \textbf{What to point at} \\
\midrule
\textbf{Overview} & Baseline integrity (\code{VALID}/\code{TAMPERED}), ``program in sync?'',
  the severity count strip, live event total. The whole-health glance. \\
\textbf{Program Diff} & The GitHub-style red/green side-by-side of baseline vs running L5X,
  word-level highlights per changed rung, with severity + MITRE per change. \emph{Use this
  for the logic mutations (3a--3e).} \\
\textbf{Alerts} & The severity-ranked feed; each alert shows its \code{who / when / channel}
  identity and its ATT\&CK-for-ICS technique. \emph{Use this for every act.} \\
\textbf{Live Plant} & The thermal-plant mimic from live Modbus reads (drum level, steam
  pressure, turbine speed, trip coils). \emph{Shows the physical effect of 3a / 4a.} \\
\textbf{Evidence} & The forensic log + one-click signed PDF report. \\
\bottomrule
\end{tabularx}

% ==================================================================
\section{Judge Q\&A}
\begin{itemize}[leftmargin=1.3em,itemsep=3pt]
  \item \textbf{``Is the PLC program real?''} Yes --- a genuine Rockwell \code{.L5X} export.
        We never \emph{execute} it; we parse, canonicalize (strip volatile
        dates/revisions/CRC/comments), and diff it. That's why a benign re-export never
        false-alarms but a single inverted instruction is caught.
  \item \textbf{``How do you know the baseline itself wasn't tampered?''} It's HMAC-SHA256
        signed; editing the approved copy breaks the signature and raises a \crit\
        \code{baseline\_tamper} on its own --- and a passive file-integrity monitor watches
        it in real time.
  \item \textbf{``Real MITRE or made up?''} Rule-based, explainable, and \emph{verified}
        against the live ATT\&CK for ICS matrix. Where no technique fits (physical enclosure
        tamper) we say \code{N/A} rather than invent an ID.
  \item \textbf{``Does it need the Pi?''} No --- it runs entirely on one laptop with an
        embedded plant for judging, and splits cleanly onto a real Pi + attacker box over
        Wi-Fi with the same code.
  \item \textbf{``How is it verified?''} 98 automated checks across 7 suites
        (\code{python -m logicward.tests.smoke\_*}).
\end{itemize}

% ==================================================================
\clearpage
\begin{landscape}
\thispagestyle{fancy}
\section{Master map --- attack $\rightarrow$ detection $\rightarrow$ MITRE}
\renewcommand{\arraystretch}{1.3}
\footnotesize
\begin{tabularx}{\linewidth}{@{}p{2.5cm} p{3.9cm} p{3.0cm} L p{2.9cm} p{3.2cm} c@{}}
\toprule
\textbf{Attack} & \textbf{Command (attacker)} & \textbf{Lands on / channel} &
\textbf{Detected as (plane)} & \textbf{Dashboard view} & \textbf{MITRE ATT\&CK for ICS} & \textbf{Sev.} \\
\midrule
Rogue device & \code{attacks rogue} & OT LAN / ARP announce &
  \code{physical.rogue\_device} (physical) & Alerts & Initial Access \textbullet\ Rogue Master \textbullet\ \code{T0848} & \med \\
Setpoint drift & \code{attacks setpoint-drift} & PLC Modbus HR / FC06 write &
  \code{cyber.setpoint\_drift} (cyber) & Alerts, Program Diff, Live Plant & Impair Process Control \textbullet\ Modify Parameter \textbullet\ \code{T0836} & \high \\
Logic inversion & \code{attacks logic-inversion} & PLC program / download &
  \code{cyber.logic\_inversion} (cyber) & Program Diff, Alerts & Persistence \textbullet\ Modify Program \textbullet\ \code{T0889} & \crit \\
Condition stripping & \code{attacks condition-stripping} & PLC program / download &
  \code{cyber.condition\_stripping} (cyber) & Program Diff, Alerts & Persistence \textbullet\ Modify Program \textbullet\ \code{T0889} & \crit \\
Coil hijack & \code{attacks coil-hijack} & PLC program / download &
  \code{cyber.coil\_hijack} (cyber) & Program Diff, Alerts & Persistence \textbullet\ Modify Program \textbullet\ \code{T0889} & \crit \\
Rung injection & \code{attacks rung-injection} & PLC program / download &
  \code{cyber.rung\_injection} (cyber) & Program Diff, Alerts & Lateral Movement \textbullet\ Program Download \textbullet\ \code{T0843} & \crit \\
(any program download) & --- (side effect of 3b--3e) & PLC program file &
  \code{cyber.program\_file\_modified} (FIM) & Alerts & Persistence \textbullet\ Modify Program \textbullet\ \code{T0889} & \low \\
Force coil & \code{attacks force-coil} & PLC Modbus coil / FC05 write &
  \code{cyber.register\_change} (cyber) & Alerts, Live Plant & Impair Process Control \textbullet\ Unauthorized Command Message \textbullet\ \code{T0855} & \low \\
DDoS flood & \code{attacks ddos --count 800} & PLC Modbus server / flood &
  \code{resource.cpu\_spike} (resource) & Alerts & Inhibit Response Function \textbullet\ Denial of Service \textbullet\ \code{T0814} & \med \\
Cable pull & (physical / \code{link\_watch}) & PLC NIC carrier &
  \code{physical.link\_down} (physical) & Alerts & Inhibit Response Function \textbullet\ Denial of Service \textbullet\ \code{T0814} & \med \\
Enclosure open & (physical / \code{gpio\_tamper}) & PLC cabinet switch &
  \code{physical.enclosure\_open} (physical) & Alerts & Initial Access \textbullet\ (no direct technique) \textbullet\ \code{N/A} & \high \\
\bottomrule
\end{tabularx}

\vspace{6pt}
\footnotesize\color{muted}
Severity is computed per event type (base weight $\times$ a safety-critical multiplier), not
hand-set. A full \code{demo\_sequence} run = 14 events (4 critical, 2 high, 3 medium, 5 low).
\end{landscape}

\end{document}
